% Propietario conservado: /Users/ruben/Documents/New project/output/EXTREMA_MEDIA_RAZON_RECIPROCIDAD_ES_EN_20260918/ARTICULO/ES/source/libro/censo_firmas_emision.tex
\chapter{Census of residual-emitter signatures}
\label{x_reciprocidad:app:firmas-emision}

This census specifies the finite fibers used in Section~\ref{x_reciprocidad:sec:extension}.
Each initial condition is identified by
\[
 \nu=4(9i+j)+d,\qquad 0\leq i,j\leq8,\quad
 d=0,1,2,3\ \longleftrightarrow\ N,E,S,O.
\]
Euclidean division recovers \(d=[\nu]_4\), \(j=[\lfloor\nu/4\rfloor]_9\)
and \(i=\lfloor\nu/36\rfloor\). Thus \(0\leq\nu<324\) enumerates
each condition exactly once, and pairs of these identifiers enumerate
the \(104\,976\) conditions in \(\Omega_{\rm seed}\).

Let \(f_+(\nu)\) and \(f_\times(\nu)\) be the six-component signatures
obtained from the emission recurrence. At every active step the cell is read
before the cursor advances; the direction vectors are
\[
 v_N=(-1,0),\quad v_E=(0,1),\quad v_S=(1,0),\quad v_O=(0,-1).
\]
Coordinates are reduced modulo nine and the vector is reversed after
steps \(27\) and \(54\). Zero-phase steps do not move either cursor.
Together with \eqref{x_reciprocidad:eq:emisor-seis}, these rules determine
\(f_+\) and \(f_\times\) without a target-value catalog.

The table gives their complete images. The multiplicity \(m_\sigma\)
counts identifiers in the fiber \(f_\sigma^{-1}(a)\); \(\nu_{\min}\)
is its least element. Each word retains its initial zeros. The complete
fiber is determined by
\[
 f_\sigma^{-1}(a)=\{\nu\in\{0,\ldots,323\}:f_\sigma(\nu)=a\}.
\]
\begin{center}\small
\begin{tabular}{@{}lrr@{\qquad}lrr@{}}
\toprule
\(f_+\) & \(m_+\) & \(\nu_{\min}\) & \(f_\times\) & \(m_\times\) & \(\nu_{\min}\)\\
\midrule
\texttt{122564} & 18 & 17 & \texttt{000000} & 108 & 8 \\
\texttt{122898} & 18 & 24 & \texttt{177393} & 8 & 1 \\
\texttt{212645} & 18 & 5 & \texttt{177555} & 8 & 54 \\
\texttt{212988} & 18 & 0 & \texttt{177771} & 8 & 11 \\
\texttt{221456} & 18 & 29 & \texttt{339555} & 8 & 6 \\
\texttt{221889} & 18 & 12 & \texttt{339771} & 8 & 15 \\
\texttt{456221} & 18 & 28 & \texttt{339933} & 8 & 59 \\
\texttt{465898} & 18 & 21 & \texttt{393177} & 8 & 0 \\
\texttt{546988} & 18 & 9 & \texttt{393393} & 8 & 45 \\
\texttt{564122} & 18 & 16 & \texttt{393555} & 8 & 40 \\
\texttt{645212} & 18 & 4 & \texttt{555177} & 8 & 52 \\
\texttt{654889} & 18 & 33 & \texttt{555339} & 8 & 4 \\
\texttt{889221} & 18 & 13 & \texttt{555393} & 8 & 41 \\
\texttt{889654} & 18 & 32 & \texttt{555555} & 24 & 84 \\
\texttt{898122} & 18 & 25 & \texttt{555717} & 8 & 14 \\
\texttt{898465} & 18 & 20 & \texttt{555771} & 8 & 18 \\
\texttt{988212} & 18 & 1 & \texttt{555933} & 8 & 24 \\
\texttt{988546} & 18 & 8 & \texttt{717555} & 8 & 12 \\
 & &  & \texttt{717717} & 8 & 21 \\
 & &  & \texttt{717933} & 8 & 19 \\
 & &  & \texttt{771177} & 8 & 9 \\
 & &  & \texttt{771339} & 8 & 13 \\
 & &  & \texttt{771555} & 8 & 16 \\
 & &  & \texttt{933339} & 8 & 57 \\
 & &  & \texttt{933555} & 8 & 26 \\
 & &  & \texttt{933717} & 8 & 17 \\
\bottomrule
\end{tabular}
\end{center}

Evaluating the recurrence for the \(324\) identifiers gives the two
partitions in the table: \(18\cdot18=324\) and
\(24\cdot8+24+108=324\). Their fiber product yields \(432\) fibers of
size \(144\), \(18\) of size \(432\), and \(18\) of size \(1944\),
by injectivity of \eqref{x_reciprocidad:eq:acoplamiento}. Finally, applying
\eqref{x_reciprocidad:eq:psi-catalogo} to that product of images gives \(243\)
distinct words. This finite census belongs to the residual emitter;
subsequent coinductive refinement retains its own law.

